% Part: computability
% Chapter: recursive-functions
% Section: bounded-minimization

\documentclass[../../../include/open-logic-section]{subfiles}

\begin{document}

\olfileid[pt]{cmp}{rec}{bmi}
\olsection{Minimização limitada}

\begin{explain}
Muitas vezes é útil definir uma função como o menor número que satisfaz certa
propriedade ou relação~$P$. Se $P$ é decidível, podemos calcular essa função
simplesmente testando todos os números possíveis, $0$, $1$, $2$, \dots, até
encontrarmos o menor que satisfaz~$P$. Esse tipo de busca ilimitada nos leva
para fora do âmbito das funções recursivas primitivas. Contudo, se estamos
interessados apenas no menor número \emph{menor que algum limite dado
independentemente}, permanecemos no âmbito recursivo primitivo. Em outras
palavras, e de modo um pouco mais geral, suponha que tenhamos uma relação
recursiva primitiva~$R(x,z)$. Considere a função que leva $x$ e~$y$ ao menor
$z < y$ tal que $R(x, z)$. Ela também pode ser calculada, testando-se
$R(x, 0)$, $R(x, 1)$, \dots, $R(x, y-1)$. Mas por que ela é recursiva
primitiva?
\end{explain}

\begin{prop}
Se $R(\vec x, z)$ é recursiva primitiva, também o é a função
$m_R(\vec{x}, y)$ que devolve o menor~$z$ menor que~$y$ tal que
$R(\vec x, z)$ vale, se houver algum, e devolve $y$ caso contrário.
Escreveremos a função~$m_R$ como
\[
\bmin{z < y}{R(\vec{x}, z)},
\]
\end{prop}

\begin{proof}
Observe que não pode haver~$z < 0$ tal que $R(\vec{x}, z)$, pois não há
nenhum $z < 0$. Logo, $m_R(\vec x, 0) = 0$.

Quando o limite tem a forma $y + 1$, temos três casos:
\begin{enumerate}
\item Existe $z < y$ tal que $R(\vec{x}, z)$; nesse caso,
$m_R(\vec{x}, y+1) = m_R(\vec{x}, y)$.
\item Não existe tal~$z<y$, mas $R(\vec{x}, y)$ vale; então,
$m_R(\vec{x}, y+1) = y$.
\item Não existe $z < y+1$ tal que $R(\vec{x}, z)$; então,
$m_R(\vec{z}, y+1) = y+1$.
\end{enumerate}
Assim, podemos definir $m_R(\vec x, 0)$ por recursão primitiva como segue:
\begin{align*}
m_R(\vec{x}, 0) & = 0\\
m_R(\vec{x}, y+1) & =
\begin{cases}
m_R(\vec{x}, y) & \text{se $m_R(\vec x, y) \neq y$}\\
y & \text{se $m_R(\vec x, y) = y$ e $R(\vec{x}, y)$}\\
y+1 & \text{caso contrário.}
\end{cases}
\end{align*}
Observe que existe $z<y$ tal que $R(\vec x, z)$ se, e somente se,
$m_R(\vec x, y) \neq y$.
\end{proof}

\begin{prob}
Suponha que $R(\vec x, z)$ seja recursiva primitiva. Defina a função
$m'_R(\vec{x}, y)$ que devolve o menor~$z$ menor que~$y$ tal que
$R(\vec{x}, z)$ vale, se houver algum, e devolve $0$ caso contrário, por
recursão primitiva a partir de~$\Char{R}$.
\end{prob}

\end{document}
